\chapter{超音速牛奶喷射与时间旅行的应用前景}

本章讨论模型从实验分析走向应用时需要补充的工程环节。应用前景不能只写成愿景，还应明确输入、处理、决策、异常处理和输出之间的对应关系，并说明哪些结论仍然依赖实验条件。

\section{应用场景}

\subsection{食品加工}

在食品加工场景中，喷射速度、温度和黏度共同决定输送稳定性。实际部署时，控制器需要根据在线传感器估计当前状态，再决定是否调整喷嘴压力或暂停生产。因此，论文应同时报告稳态性能和扰动恢复时间，而不应只报告理想工况下的峰值效率。

\subsection{时空导航}

时空导航场景更关心状态预测和不确定性传播。模型输出不能只给出一个时间点，而应给出预测区间、置信水平和异常状态的处理策略。对于跨设备部署，还要明确时间同步误差、通信延迟和模型更新周期。

\subsection{数据处理算法}

论文描述数值求解、实验处理或数据筛选过程时，可以把关键步骤整理为带缩进和行号的伪代码。算法~\ref{alg:response}与前文的状态模型对应，强调执行顺序而不是某一种编程语言的语法。
\begin{algorithm}[htbp]
  \caption{概率响应估计算法}
  \label{alg:response}
  \begin{algorithmic}[1]
    \REQUIRE 观测序列 $\bm{y}_1,\ldots,\bm{y}_n$，参数 $\theta$，异常阈值 $\epsilon$
    \ENSURE 估计状态 $\hat{\bm{x}}_1,\ldots,\hat{\bm{x}}_n$ 与异常标记
    \STATE 读取观测序列与参数 $\theta$
    \STATE 初始化状态 $\bm{x}_0$ 和协方差 $\bm{P}_0$
    \FOR{$k=1,2,\ldots,n$}
      \STATE 预测状态并计算残差
      \STATE 更新状态与置信区间
      \IF{残差超过阈值 $\epsilon$}
        \STATE 标记异常观测，并记录观测序号
      \ENDIF
    \ENDFOR
    \STATE 输出状态序列与异常记录
  \end{algorithmic}
\end{algorithm}

算法~\ref{alg:response}中的阈值、状态转移矩阵和协方差更新公式应在正文中分别定义。这样读者能够把伪代码中的每一步追溯到模型假设，而不会把算法框误认为不可复现的流程装饰。

\subsection{流程示意图}

伪代码强调步骤顺序，流程图则适合说明分支和反馈关系。图~\ref{fig:process}给出数据读取、参数估计、异常重试和结果输出之间的关系。
\begin{figure}[htbp]
  \centering
  \begin{tikzpicture}[>=Stealth,node distance=15mm,
    block/.style={draw,rounded corners,align=center,minimum width=27mm,minimum height=10mm}]
    \node[block,fill=blue!7] (read) {读取数据};
    \node[block,fill=teal!7,right=of read] (estimate) {估计参数};
    \node[block,fill=yellow!12,right=of estimate] (output) {输出结果};
    \node[block,fill=orange!10,below=13mm of estimate] (retry) {调整步长};
    \draw[->] (read)--(estimate);
    \draw[->] (estimate)--(output);
    \draw[->] (estimate)--(retry);
    \draw[->] (retry.west)-| (read.south);
  \end{tikzpicture}
  \caption{带异常重试的参数估计流程}
  \label{fig:process}
\end{figure}

\subsection{模块接口关系}

当论文需要强调模块接口而不是几何位置时，可以用紧凑的节点和箭头表达状态转移。图~\ref{fig:module-interface}把数据采集、估计和决策模块按执行顺序连接起来。示意图应控制在版心宽度以内，并通过正文解释节点含义，不能把图例文字全部塞进图内。
\begin{figure}[htbp]
  \centering
  \begin{tikzpicture}[>=Stealth,node distance=18mm,
    module/.style={draw,rounded corners,fill=gray!10,minimum width=25mm,minimum height=10mm,align=center}]
    \node[module] (collect) {采集};
    \node[module,right=of collect] (estimate) {估计};
    \node[module,right=of estimate] (decision) {决策};
    \draw[->,thick] (collect)--(estimate);
    \draw[->,thick] (estimate)--(decision);
  \end{tikzpicture}
  \caption{状态处理模块的接口关系}
  \label{fig:module-interface}
\end{figure}

图~\ref{fig:engineering-diagrams}把六类工程结构并列为一组：层级架构、输入输出流程、有向网络、节点分组、UML 接口和代数结构。总题注说明共同主题，子图题注分别说明各自表达的关系；与本章论证关系不大的图可保留为独立普通图。
\begin{figure}[htbp]
  \centering
  \begin{subfigure}[t]{0.49\textwidth}
    \centering
    \resizebox{\linewidth}{!}{%
      \begin{tikzpicture}[>=Stealth, every node/.style={font=\scriptsize},
        layer/.style={draw,rounded corners=3pt,align=center,minimum height=7mm,inner sep=3pt},
        norm/.style={layer,shape=diamond,aspect=2.8,fill=yellow!35,draw=orange!80!black,inner sep=1pt},
        inputnode/.style={layer,shape=trapezium,trapezium angle=70,fill=red!20,draw=red!75!black}]
        \useasboundingbox (-0.1,-0.1) rectangle (7.2,3.5);
        \node[inputnode,minimum width=10mm] (input) at (1.45,0.25) {Input};
        \node[layer,fill=orange!30,draw=orange!80!black,text width=19mm,inner sep=1pt] (attention) at (1.45,1.25) {Multi-Head\\Attention};
        \node[norm,minimum width=18mm] (norm-a) at (1.45,2.35) {Add \& Norm};
        \node[layer,fill=blue!25,draw=blue!75!black,text width=18mm,inner sep=1pt] (feed) at (4.85,1.25) {Feed\\Forward};
        \node[norm,minimum width=18mm] (norm-b) at (4.85,2.35) {Add \& Norm};
        \node[layer,fill=green!25,draw=green!65!black,minimum width=15mm] (softmax) at (6.55,2.35) {Softmax};
        \draw[->,thick] (input.north) -- (attention.south);
        \draw[->,thick] (attention.north) -- (norm-a.south);
        \draw[->,thick,rounded corners=2pt] (norm-a.east) -- (3.25,2.35) -- (3.25,1.25) -- (feed.west);
        \draw[->,thick] (feed.north) -- (norm-b.south);
        \draw[->,thick] (norm-b.east) -- (softmax.west);
        \draw[->,dashed,rounded corners=3pt] (input.west) -- (0.3,0.25) -- (0.3,2.35) -- (norm-a.west);
        \draw[->,dashed,rounded corners=3pt] (norm-a.north) -- (1.45,3.12) -- (4.85,3.12) -- (norm-b.north);
      \end{tikzpicture}%
    }
    \caption{层级架构}
  \end{subfigure}\hfill
  \begin{subfigure}[t]{0.49\textwidth}
    \centering
    \resizebox{\linewidth}{!}{%
      \begin{tikzpicture}[>=Stealth, every node/.style={font=\scriptsize},
        unit/.style={draw,rounded corners=4pt,align=center,minimum width=25mm,minimum height=8mm,fill=green!35,draw=green!55!black},
        port/.style={draw,align=center,minimum width=14mm,minimum height=7mm}]
        \useasboundingbox (-0.1,-0.1) rectangle (7.2,3.5);
        \node[port,fill=magenta!18,draw=magenta!60!black] (input) at (0.72,1.55) {input};
        \node[unit] (cu) at (3.45,2.85) {control unit (CU)};
        \node[unit,text width=31mm] (alu) at (3.45,1.5) {arithmetic \& logic\\unit (ALU)};
        \node[unit] (mu) at (3.45,0.15) {memory unit (MU)};
        \node[port,shape=regular polygon,regular polygon sides=6,fill=cyan!25,draw=cyan!55!black] (output) at (6.25,1.5) {output};
        \draw[->,thick,rounded corners=4pt] (input.east) -- (1.75,1.5) -- (1.75,2.85) -- (cu.west);
        \draw[<->,thick] (cu.south) -- (alu.north);
        \draw[<->,thick] (alu.south) -- (mu.north);
        \draw[->,thick,rounded corners=4pt] (cu.east) -- (5.0,2.85) -- (5.0,1.5) -- (output.west);
        \draw[->,dashed,rounded corners=4pt] (mu.east) -- (5.35,0.15) -- (5.35,2.45) -- (cu.east);
      \end{tikzpicture}%
    }
    \caption{输入输出流程}
  \end{subfigure}

  \par\smallskip

  \begin{subfigure}[t]{0.49\textwidth}
    \centering
    \resizebox{\linewidth}{!}{%
      \begin{tikzpicture}[>=Stealth, every node/.style={font=\small},
        vertex/.style={draw,circle,minimum size=8mm,inner sep=1pt,fill=white}]
        \useasboundingbox (-0.1,-0.1) rectangle (7.2,3.5);
        \node[vertex] (a) at (3.55,3.0) {A};
        \node[vertex] (b) at (5.35,2.28) {B};
        \node[vertex] (c) at (5.95,1.13) {C};
        \node[vertex] (d) at (4.75,0.35) {D};
        \node[vertex] (e) at (2.35,0.35) {E};
        \node[vertex] (f) at (1.05,1.13) {F};
        \node[vertex] (g) at (1.62,2.28) {G};
        \draw[->] (d) -- (c);
        \draw[->] (e) to[bend left=18] (b);
        \draw[->] (b) to[bend left=18] (e);
        \draw[->] (a) -- (e);
        \draw[->] (d) -- (a);
        \draw[->] (f) to[bend left=20] (g);
        \draw[->] (g) to[bend left=20] (f);
      \end{tikzpicture}%
    }
    \caption{有向网络}
  \end{subfigure}\hfill
  \begin{subfigure}[t]{0.49\textwidth}
    \centering
    \resizebox{\linewidth}{!}{%
      \begin{tikzpicture}[>=Stealth, every node/.style={font=\scriptsize},
        obj/.style={font=\small}]
        \useasboundingbox (-0.1,-0.1) rectangle (7.2,3.5);
        \draw[rounded corners=4pt,draw=cyan!70!black,fill=cyan!4,line width=1pt] (0.15,1.2) rectangle (3.55,3.25);
        \draw[rounded corners=4pt,draw=cyan!70!black,fill=cyan!4,line width=1pt] (4.25,1.2) rectangle (6.95,3.25);
        \draw[rounded corners=4pt,draw=green!65!black,fill=green!4,line width=1pt] (0.15,0.08) rectangle (3.55,0.95);
        \draw[rounded corners=4pt,draw=green!65!black,fill=green!4,line width=1pt] (4.25,0.08) rectangle (6.95,0.95);
        \node[obj] (sa) at (1.0,2.78) {$Sa$};
        \node[obj] (tb) at (2.75,2.78) {$Tb$};
        \node[obj] (sap) at (1.0,1.65) {$Sa'$};
        \node[obj] (tbp) at (2.75,1.65) {$Tb'$};
        \node[obj] (tuple) at (5.6,2.78) {$(a,b,f)$};
        \node[obj] (tuplep) at (5.6,1.65) {$(a',b',f')$};
        \node[obj] (sa-small) at (1.0,0.52) {$Sa$};
        \node[obj] (tb-small) at (2.75,0.52) {$Tb$};
        \node[obj] (tuple-small) at (5.6,0.52) {$(a,b,f)$};
        \draw[->] (sa) -- node[above] {$f$} (tb);
        \draw[->] (sap) -- node[above] {$f'$} (tbp);
        \draw[->] (sa) -- node[left] {$\alpha$} (sap);
        \draw[->] (tb) -- node[right] {$\beta$} (tbp);
        \draw[->] (tuple) -- node[right] {$(\alpha,\beta)$} (tuplep);
        \draw[<->,cyan!70!black,line width=1pt] (3.62,2.2) -- (4.18,2.2);
        \draw[<->,black,line width=0.8pt] (1.85,1.08) -- (1.85,1.17);
        \draw[<->,black,line width=0.8pt] (5.6,1.08) -- (5.6,1.17);
        \draw[->] (sa-small) -- node[above] {$f$} (tb-small);
        \draw[<->,green!65!black,line width=1pt] (3.62,0.52) -- (4.18,0.52);
      \end{tikzpicture}%
    }
    \caption{节点分组}
  \end{subfigure}

  \par\smallskip

  \begin{subfigure}[t]{0.49\textwidth}
    \centering
    \resizebox{\linewidth}{!}{%
      \begin{tikzpicture}[>=Stealth, every node/.style={font=\scriptsize},
        class/.style={draw=gray!65,fill=white,align=center,minimum width=19mm,minimum height=8mm,font=\bfseries}]
        \useasboundingbox (-0.1,-0.1) rectangle (7.2,3.5);
        \node[class] (diagram) at (3.55,2.92) {Diagram};
        \node[class] (node) at (1.45,1.45) {Node};
        \node[class] (edge) at (5.65,1.45) {Edge};
        \node[class] (mark) at (5.65,0.3) {Mark};
        \draw (diagram.south) -- (3.55,2.35) -- (1.45,2.35) -- (node.north);
        \draw (3.55,2.35) -- (5.65,2.35) -- (edge.north);
        \node[above left,inner sep=1pt] at (3.48,2.87) {$1$};
        \node[above,inner sep=1pt] at (1.45,1.95) {$0..n$};
        \node[above,inner sep=1pt] at (5.65,1.95) {$0..n$};
        \draw (edge.south) -- node[right,fill=white,inner sep=1pt] {$1..n$} (mark.north);
        \draw[->,cyan!70!black] (edge.west) -- node[above,fill=white,inner sep=1pt] {snap} (node.east);
        \draw[->,orange!85!black,bend left=34] (node.west) to node[left,fill=white,inner sep=1pt] {layout} (diagram.west);
      \end{tikzpicture}%
    }
    \caption{UML 接口关系}
  \end{subfigure}\hfill
  \begin{subfigure}[t]{0.49\textwidth}
    \centering
    \resizebox{\linewidth}{!}{%
      \begin{tikzpicture}[>=Stealth, every node/.style={font=\scriptsize},
        algebra/.style={draw=none,fill=gray!8,inner sep=2pt,minimum height=5mm,align=center}]
        \useasboundingbox (-0.1,-0.1) rectangle (7.2,3.5);
        \node[algebra] (magma) at (3.6,3.2) {magma};
        \node[algebra] (semi) at (1.25,2.35) {semigroup};
        \node[algebra] (unital) at (3.6,2.35) {unital magma};
        \node[algebra] (quasi) at (5.95,2.35) {quasigroup};
        \node[algebra] (monoid) at (1.25,1.35) {monoid};
        \node[algebra] (invsemi) at (3.6,1.35) {inverse semigroup};
        \node[algebra] (loop) at (5.95,1.35) {loop};
        \node[algebra] (group) at (3.6,0.35) {group};
        \draw[->,blue!75!black] (magma) -- node[above left,fill=white,inner sep=1pt] {Assoc} (semi);
        \draw[->,red!75!black] (magma) -- node[right,fill=white,inner sep=1pt] {Id} (unital);
        \draw[->,yellow!65!black] (magma) -- node[above right,fill=white,inner sep=1pt] {Div} (quasi);
        \draw[->,red!75!black] (semi) -- node[left,fill=white,inner sep=1pt] {Id} (monoid);
        \draw[->,yellow!65!black] (semi) to[bend right=8] (invsemi);
        \node[font=\scriptsize,text=yellow!65!black,fill=white,inner sep=1pt] at (2.0,2.08) {Div};
        \draw[->,blue!75!black] (unital) to[bend right=8] (monoid);
        \node[font=\scriptsize,text=blue!75!black,fill=white,inner sep=1pt] at (2.85,1.62) {Assoc};
        \draw[->,green!60!black] (unital) -- node[right,fill=white,inner sep=1pt] {Inv} (invsemi);
        \draw[->,blue!75!black] (quasi) to[bend left=8] node[above,fill=white,inner sep=1pt] {Assoc} (invsemi);
        \draw[->,red!75!black] (quasi) -- node[right,fill=white,inner sep=1pt] {Id} (loop);
        \draw[->,green!60!black] (monoid) -- node[above left,fill=white,inner sep=1pt] {Inv} (group);
        \draw[->,red!75!black] (invsemi) -- node[right,fill=white,inner sep=1pt] {Id} (group);
        \draw[->,blue!75!black] (loop) -- node[above right,fill=white,inner sep=1pt] {Assoc} (group);
      \end{tikzpicture}%
    }
    \caption{代数结构层级}
  \end{subfigure}
  \caption{工程模块图组：层级架构、输入输出流程、有向网络、节点分组、UML 接口与代数结构}
  \label{fig:engineering-diagrams}
\end{figure}

\section{数据处理与评估}

\subsection{滤波与谱分析}

实测信号总含有测量噪声，因此进入参数估计之前需要先做滤波。滑动平均是最简单的低通滤波，一阶惯性环节则给出可解析的传递函数，二者都应在正文中说明截止频率与相位滞后。
\begin{equation}
  \hat{y}_k=\frac{1}{2m+1}\sum_{j=-m}^{m}y_{k+j},\qquad
  H(\omega)=\frac{1}{1+\mathrm{i}\omega\tau},\qquad
  |H(\omega)|=\frac{1}{\sqrt{1+(\omega\tau)^2}}.
  \label{eq:filters}
\end{equation}

功率谱密度把信号的方差按频率分解，用于判断周期性扰动来自哪一频段。估计谱时需要说明窗函数和平均段数，否则峰值会被泄漏和方差掩盖。
\begin{equation}
  S(\omega)=\lim_{T\to\infty}\frac{1}{T}
  \left|\int_0^T y(t)e^{-\mathrm{i}\omega t}\,\mathrm{d}t\right|^2,\qquad
  \int_{-\infty}^{\infty}S(\omega)\,\mathrm{d}\omega=\operatorname{Var}[y].
  \label{eq:psd}
\end{equation}

离散频谱的可用范围由采样定理限定，混叠一旦发生就无法通过后处理恢复。工程上通常在采样前加模拟低通滤波，并在正文中给出抗混叠滤波器的参数。
\begin{equation}
  \omega_{\mathrm{Nyquist}}=\pi f_s,\qquad
  f_s=\frac{1}{\Delta t},\qquad f_s>2f_{\max}.
  \label{eq:nyquist}
\end{equation}

\subsection{指标与不确定度}

测量结果的不确定度由各输入量的不确定度按灵敏度系数合成。输入量之间相关时必须保留协方差项，只有相互独立时才能把平方和简化成对角线形式。
\begin{equation}
  u_c^2(y)=
  \sum_{i=1}^{n}\left(\frac{\partial f}{\partial x_i}\right)^2u^2(x_i)
  +2\sum_{i<j}\frac{\partial f}{\partial x_i}
  \frac{\partial f}{\partial x_j}r_{ij}u(x_i)u(x_j).
  \label{eq:uncertainty-propagation}
\end{equation}
\begin{equation}
  u_c(y)=\sqrt{\sum_{i=1}^{n}
  \left(\frac{\partial f}{\partial x_i}\right)^2u^2(x_i)},\qquad
  U=ku_c,\qquad k=2,\qquad P\approx95\%.
  \label{eq:expanded-uncertainty}
\end{equation}

评估指标必须成组报告：单看平均误差会掩盖大偏差样本，单看最大误差又会被个别异常值主导。表~\ref{tab:metrics}给出示例中使用的指标及其适用条件。
\begin{table}[htbp]
  \centering
  \caption{评估指标的定义与适用条件}
  \label{tab:metrics}
  \begin{tabularx}{\textwidth}{M{0.16\textwidth} C C}
    \toprule
    \textbf{指标} & \textbf{定义} & \textbf{适用条件}\\
    \midrule
    RMSE & 误差平方均值的平方根 & 重视大偏差、误差近似对称\\
    MAE & 绝对误差的平均值 & 存在异常值、需要稳健性\\
    $R^2$ & 残差平方和与总平方和之比 & 需要与基线模型比较\\
    覆盖率 & 真值落入区间的比例 & 报告区间估计的可靠性\\
    \bottomrule
  \end{tabularx}
\end{table}

偏差与精密度是两类不同的误差。前者反映系统性偏离，可以通过标定消除；后者反映随机分散，只能通过增加样本量或改进测量链降低。
\begin{equation}
  \mathrm{bias}=\frac{1}{n}\sum_{i=1}^{n}(\hat{y}_i-y_i),\qquad
  s=\sqrt{\frac{1}{n-1}\sum_{i=1}^{n}
  \left(\hat{y}_i-\overline{\hat{y}}\right)^2}.
  \label{eq:bias-precision}
\end{equation}

\subsection{带单位表头与表注的表}

单位写在表头里比写在每个数据后面更紧凑，也便于整列换算；表注放在表体下方、与表同属一个浮动体，用于说明数据来源、符号含义和保留位数。
\begin{table}[htbp]
  \centering
  \caption{带单位表头与表注的参数表}
  \label{tab:units}
  \begin{tabular}{ccccc}
    \toprule
    \textbf{工况} & $v_0$/(\si{\meter\per\second}) & $p_0$/\si{\kilo\pascal} & $T_0$/\si{\kelvin} & 样本量 $n$\\
    \midrule
    工况 A & 340.0 & 101.3 & 298.1 & 12\\
    工况 B & 352.5 & 104.7 & 301.4 & 12\\
    工况 C & 361.2 & 98.6 & 296.3 & 15\\
    工况 D & 375.8 & 107.2 & 303.9 & 15\\
    \bottomrule
  \end{tabular}

  \smallskip
  \begin{minipage}{0.9\textwidth}\footnotesize
    注：表中为示例数据，$v_0$ 为喷嘴出口速度，$p_0$ 与 $T_0$ 为滞止参数；保留位数按仪器分辨力给出。
  \end{minipage}
\end{table}

\subsection{并排小表}

两组独立的对照数据可以并排放在一个版面中，各自保留表号、表题和三线。并排时应保证合计宽度不超过版心；表格内容较多时不要并排，以免单个单元格过窄而频繁换行。
\begin{center}
  \begin{minipage}[t]{0.48\textwidth}
    \centering
    \captionsetup{type=table}
    \caption{两组工况的误差指标}
    \label{tab:metrics-ab}
    \begin{tabular}{ccc}
      \toprule
      指标 & 工况 A & 工况 B\\
      \midrule
      RMSE & 0.0241 & 0.0198\\
      MAE & 0.0187 & 0.0152\\
      $R^2$ & 0.961 & 0.973\\
      \bottomrule
    \end{tabular}
  \end{minipage}
  \hfill
  \begin{minipage}[t]{0.48\textwidth}
    \centering
    \captionsetup{type=table}
    \caption{另外两组工况的误差指标}
    \label{tab:metrics-cd}
    \begin{tabular}{ccc}
      \toprule
      指标 & 工况 C & 工况 D\\
      \midrule
      RMSE & 0.0212 & 0.0175\\
      MAE & 0.0163 & 0.0138\\
      $R^2$ & 0.968 & 0.979\\
      \bottomrule
    \end{tabular}
  \end{minipage}
\end{center}

\subsection{计算复杂度与资源}

把模型部署到在线设备上时，需要同时说明时间复杂度和存储需求。下式给出按网格数估计的计算量与吞吐率关系，正文中应给出实测值与估计值的差异，而不应只给阶数。
\begin{equation}
  T(n)=Cn^3,\qquad M(n)=4n^2\ \text{字节},\qquad
  t_{\mathrm{step}}=\frac{T(n)}{\text{throughput}},\qquad
  E=Pt_{\mathrm{step}}.
  \label{eq:complexity}
\end{equation}

迭代求解的终止条件也需要量化。常用的判据是相对残差或参数增量小于给定阈值，阈值应与测量不确定度相当，过严的阈值只会增加迭代次数而不改善结果。
\begin{equation}
  \frac{\|\bm{r}^{(k)}\|}{\|\bm{b}\|}<\epsilon_{\mathrm{rel}},\qquad
  \frac{\|\bm{\theta}^{(k+1)}-\bm{\theta}^{(k)}\|}
  {\|\bm{\theta}^{(k)}\|}<\epsilon_{\mathrm{step}}.
  \label{eq:stopping}
\end{equation}

\subsection{置信区间估计算法}

算法~\ref{alg:bootstrap}给出重采样置信区间的处理流程。该方法不假设统计量的解析分布，因此在指标分布未知时比正态近似更稳健；代价是需要重复求解，重采样次数应按目标精度选择。
\begin{algorithm}[htbp]
  \caption{重采样置信区间估计}
  \label{alg:bootstrap}
  \begin{algorithmic}[1]
    \REQUIRE 样本 $\bm{y}_1,\ldots,\bm{y}_n$，置信水平 $1-\alpha$，重采样次数 $B$
    \ENSURE 区间估计 $[\hat{\theta}_{\mathrm{low}},\hat{\theta}_{\mathrm{high}}]$ 与相对半宽
    \STATE 计算样本统计量 $\hat{\theta}$ 与残差标准差 $\hat{\sigma}$
    \FOR{$b=1,2,\ldots,B$}
      \STATE 有放回抽取 $n$ 个样本，得到重采样样本 $\bm{y}^{(b)}$
      \STATE 在重采样样本上重新求解模型，得到 $\hat{\theta}_b$
      \IF{$\hat{\theta}_b$ 超出当前区间}
        \STATE 更新区间端点，并累计越界次数
      \ENDIF
    \ENDFOR
    \STATE 取重采样统计量的 $\alpha/2$ 与 $1-\alpha/2$ 分位数
    \IF{相对半宽大于 $\epsilon$}
      \STATE 增大 $B$ 后重新执行重采样
    \ENDIF
    \STATE 输出区间、相对半宽与越界次数
  \end{algorithmic}
\end{algorithm}

算法~\ref{alg:bootstrap}的输出包含相对半宽，便于判断重采样次数是否足够；若半宽仍大于容差，应在正文中说明是样本量不足还是模型本身对扰动敏感。

\section{风险与展望}

\subsection{工程限制}

应用研究需要同时说明设备、材料、环境和伦理等方面的约束，避免把理论推演直接表述为工程结论。模型上线前还应完成独立数据验证、异常场景测试和版本归档。

\subsection{可复现性}

建议在附录或项目仓库中保存原始数据摘要、参数文件、算法版本和编译环境。图表应由固定输入重新生成，正文中只保留经过解释的结果；这样模板替换成真实论文内容后，章节结构仍然稳定。

\section{本章小结}

本章用算法环境表达计算步骤，并用流程图说明带反馈的处理过程和模块接口。图表与伪代码共同服务于方法说明；替换为真实研究内容时，应把输入、阈值、输出和异常处理与实际实现逐项核对。
